\documentclass[11pt,a4paper]{article}
\usepackage{_preamble}



\title[Règlement intérieur -- version du \today]{Règlement intérieur}
\subtitle{Version du \today}
\author{}
\date{}

\begin{document}
   \mainHeader{}
   \maketitle

%\thispagestyle{firstpage}
%\vspace{-30pt}

   Ce règlement a comme objectif d'éviter tout accident, incident et malentendu qui pourraient nuire à l'ambiance et à la bonne tenue du Centre Équestre.

   Il s'applique dans toute l'enceinte du Centre Équestre (intérieure et extérieure).
   Il est à respecter par toute personne s'y trouvant : membres, stagiaires, propriétaires d'équidés, personnels, clientèle\dots

   Son non-respect peut entraîner une exclusion provisoire ou définitive du Centre Équestre, sans indemnité.
   Le présent règlement (accompagné de ses annexes) fait autorité sur tout autre document produit dans le cadre des activités du Centre Équestre.

   \subsection{Organisation générale}\label{subsec:organisation-generale}
   Toutes les activités du Centre Équestre ainsi que toutes les installations dont il dispose sont placées sous l'autorité des gérants ou de leur délégué.

   En période de congés ou d'absence du responsable, les employés restant sur le site sont chargés de veiller à son bon fonctionnement et d'appliquer les consignes données.

   \subsection{Fonctionnement}\label{subsec:fonctionnement}
   Pour des raisons de sécurité, tout cavalier doit attendre la présence d'un encadrant du club, ou son accord, avant de s'occuper d'un cheval (hors propriétaire s'occupant de son équidé).

   Tout cavalier doit impérativement passer au bureau avant le début de la séance pour signaler sa présence à la personne chargée de l'accueil.

   Le traitement des montoirs et cartes est informatisé.
   Ainsi, le cavalier s'engage à signaler, au plus tard la veille, son absence à une activité à laquelle il se serait inscrit, pour garantir une bonne gestion du programme et de la cavalerie.
   Par ailleurs, \emph{les leçons retenues, non décommandées avant 18h00 la veille de la reprise, sont dues}.

   Il est recommandé aux cavaliers élèves d'être présents une demi-heure avant la reprise.

   Les cartes sont strictement nominatives et ne peuvent être revendues ou transférées sans accord de l'exploitant.
   Ces dernières ne sont ni reprises ni remboursées.

   Après chaque activité, le matériel doit être remis en place (sauf avis contraire de l'enseignant) :
   \begin{itemize}
      \item
      Les selles rangées pommeau vers le mur ;
      \item
      Les tapis aux emplacements prévus ;
      \item
      Les mors lavés ;
      \item
      Les brosses, graisses et accessoires rangés aux places qui leur sont désignées.
   \end{itemize}

   Les tarifs sont affichés à l'accueil, et consultables sur le site internet.

   \subsection{Statut d'adhérent}\label{subsec:statut-d'adherent}
   Les cavaliers ont la possibilité, annuellement, de souscrire une adhésion au Centre Équestre.
   Cette dernière est valable en année scolaire (de septembre à août).
   Le montant de l'adhésion est fixé annuellement, et calculé au \textit{prorata temporis} en cas d'adhésion au-delà du 1\ier{} décembre.

   On appelle donc \enquote{adhérent} tout cavalier ayant souscrit une adhésion pour l'année en cours.

   L'adhérent a accès aux avantages fondamentaux suivants :
   \begin{itemize}
      \item tarifs préférentiels pour les cours et forfaits ;
      \item priorité d'inscription pour les stages et événements du Centre Équestre ;
      \item invitation aux activités, animations, concours internes\dots
   \end{itemize}

   À compter du 1\ier{} septembre 2026, tout cavalier adhérent a, de plus, accès aux offres suivantes :
   \begin{itemize}
      \item $15\, \%$ de réduction sur les stages\footnote{hors juillet/août} ;
      \item $15\, \%$ de réduction sur les cartes cadeau pour les balades ;
      \item $50\, \%$ de réduction pour un cours particulier pris au cours de son mois d'anniversaire ;
      \item autres avantages annuels définis en fonction de la saison.
   \end{itemize}

   \subsection{Utilisation des installations}\label{subsec:utilisation-des-installations}
   Les installations du Centre Équestre (carrières, manège, écuries\dots) ne sont utilisables que par :
   \begin{itemize}
      \item
      Les adhérents à jour de leur cotisation et de leur licence, qui participent à une activité dirigée par des encadrants du Centre Équestre ;
      \item
      Les cavaliers \enquote{de passage} qui participent à une activité encadrée par le Centre Équestre ;
      \item
      Les propriétaires qui utilisent seuls leur cheval (\textit{cf.} article~\ref{subsec:assurances}).
   \end{itemize}

   Tout autre cas d'utilisation des installations est formellement interdit.
   Une éventuelle dérogation requiert obligatoirement un accord préalable écrit du Centre Équestre.

   La pratique en autonomie du saut d'obstacles est interdite aux cavaliers non titulaires du Galop 7.

   Chaque membre doit veiller au respect des installations et à la propreté générale des lieux.

   L'accès aux écuries, pour le public, est interdit en dehors des horaires d'ouverture du Centre Équestre, afin de respecter le repos des équidés.
   Pour les propriétaires, hors cas particuliers (cheval malade, autorisation préalable de la direction\dots), l'accès aux écuries et aux installations est \textbf{interdit avant 7h30 et après 21h30}.

   Des restrictions ponctuelles ou permanentes concernant l'utilisation des installations et aires de travail peuvent être instaurées par la direction et seront notifiées aux cavaliers et propriétaires.

   Sont toujours prioritaires les activités encadrées par le Centre Équestre, leur planning étant disponible à l'accueil.

   Les véhicules doivent être garés sur le parking ou des emplacements en bordure de route prévus à cet effet.

   \subsection{Interdiction de fumer}\label{subsec:interdiction-de-fumer}

   Il est strictement interdit de fumer dans l'enceinte du Centre Équestre, à l'exception des endroits prévus à cet effet.
   Il en est de même pour la cigarette électronique.

   Les mégots devront impérativement être jetés dans un cendrier.

   \subsection{Discipline}\label{subsec:discipline}
   Durant les activités, et en particulier à l'intérieur des locaux ou installations, les adhérents doivent observer une obéissance complète aux encadrants et aux consignes données, pour la sécurité.

   En tout lieu et en toute circonstance, les adhérents sont tenus d'observer une attitude respectueuse empreinte de correction et de savoir-vivre vis-à-vis du personnel encadrant, des personnels, des autres cavaliers et des visiteurs.
   Cette attitude est évidemment à double sens.

   Aucune manifestation discourtoise envers le Centre Équestre, ses structures, ses administrateurs, son personnel ou ses adhérents n'est admise quelle qu'en soit sa forme.

   Toute personne fréquentant le Centre Équestre se doit d'observer les règles de sécurité internes (ne pas courir, faire de gestes brusques\dots) dans les écuries ou les lieux de travail.
   Les jeux (balles, ballons\dots) sont interdits.
   Il en va de même pour les jeux dans les bottes de paille ou de foin.
   Les chiens extérieurs à ceux appartenant au Centre Équestre doivent être tenus en laisse.

   D'une manière générale, les enfants ne doivent pas jouer sur le site sans être surveillés de leurs parents.
   Ils sont sous leur entière responsabilité.

   Pendant les reprises, seul l'enseignant et ses aides sont habilités à se trouver sur le terrain d'évolution.

   Il est interdit de sortir ou de déplacer les obstacles des carrières ou manège sans l'autorisation du responsable.

   \subsection{Réclamations}\label{subsec:reclamations}
   Tout adhérent souhaitant formuler une réclamation qu'il estime motivée et justifiée concernant le Centre Équestre, ses membres ou personnels, peut opérer de l'une des manières suivantes :
   \begin{itemize}
      \item
      s'adresser directement à l'un des encadrants de façon confidentielle ;
      \item
      s'adresser directement au gérant de l'exploitation par oral ou lui adresser un courrier.
   \end{itemize}
   Toute réclamation présentée sous l'une des formes ainsi définies doit recevoir une réponse sous une semaine.

   Tout adhérent est également libre de formuler des observations ou suggestions, en se rapprochant du gérant ou son délégué.

   \subsection{Sanctions}\label{subsec:sanctions}
   Toute attitude répréhensible d'un adhérent et en particulier toute inobservation du règlement intérieur expose celui qui en est responsable à des sanctions qui peuvent être de trois ordres :
   \begin{enumerate}
      \item \textbf{La mise à pied} prononcée par la direction pour une durée ne pouvant excéder un mois.
      L'adhérent qui est mis à pied ne peut, pendant la durée de la sanction, ni monter un cheval appartenant au Centre Équestre, ni utiliser les terrains d'évolution, manège ou carrières.
      \item \textbf{L'exclusion temporaire} ou suspension, prononcée par la direction pour une durée ne pouvant excéder une année.
      L'adhérent qui est exclu temporairement n'a plus accès aux locaux et installations du Centre Équestre, et ne peut, pendant la durée de la sanction, participer à aucune des activités (publique ou privée).
      Les pensions, le cas échéant, restent dues.
      \item \textbf{L'exclusion définitive}.
   \end{enumerate}
   Tout adhérent faisant l'objet d'une sanction ne peut prétendre à aucun remboursement des sommes déjà payées par lui se rapportant aux activités dont la sanction le prive.

   \emph{Le respect du cheval étant une règle absolue, tout mauvais traitement peut entraîner l'exclusion immédiate.}

   \subsection{Tenue}\label{subsec:tenue}
   Pour monter à cheval, tant à l'intérieur qu'à l'extérieur, les cavaliers doivent adopter une tenue vestimentaire correcte et conforme aux usages traditionnels de l'équitation française.

   L'enseignant est libre de refuser dans une activité qu'il encadre toute personne dont la tenue vestimentaire serait inadaptée au point de compromettre la sécurité.

   \emph{Le port du casque est obligatoire.}
   Il doit être porté afin de constituer une protection effective pour le cavalier et être conforme à la norme imposée par la règlementation en vigueur.

   \subsection{Assurances}\label{subsec:assurances}
   Le Centre Équestre est assuré pour les activités qu'il organise et encadre.
   De plus, il est rappelé aux adhérents que leur licence les assure lors de leur pratique.
   Il leur appartient de prendre connaissance au bureau de l'étendue et des limites de garantie qui leur sont ainsi accordées.

   Aucun adhérent ne peut participer aux activités du Centre Équestre ou utiliser ses installations s'il n'a pas acquitté sa cotisation pour l'année en cours et s'il n'est pas titulaire d'une licence fédérale pour l'année en cours.

   La responsabilité du Centre Équestre est dégagée dans le cas d'un accident provoqué par une inobservation du règlement intérieur.

   Les propriétaires pratiquant en autonomie sont seuls responsables.

   \subsection{Prise en charge des mineurs}\label{subsec:prise-en-charge-des-enfants-mineurs}
   Les cavaliers mineurs ne sont sous la responsabilité du Centre Équestre que durant la durée de leur reprise (ou de leur balade) et durant le temps de préparation de l'équidé et le retour à l'écurie (ou en prairie), soit généralement \textit{\textbf{une demi-heure avant}} et \textit{\textbf{une demi-heure après}}.

   En dehors des heures d'activité ainsi définies, même sur le site du Centre Équestre, \textit{\textbf{les mineurs sont sous la responsabilité de leurs parents ou de leur responsable légal}}.

   \subsection{Balades}\label{subsec:balades}

   Tout cavalier, y compris non adhérent se doit de prendre connaissance des \textit{Conditions Générales applicables aux Balades}, décrivant leurs modalités et leur fonctionnement, ainsi que les conditions de remboursement.


   \subsection{Règlement des prestations et remboursement}\label{subsec:reglement-des-prestations-et-remboursement}

   Les prestations doivent être réglées à l'accueil du Centre Équestre.
   Toute activité est conventionnellement réglée avant son début, et la direction se réserve le droit de refuser quiconque souhaite prendre part à une activité sans l'avoir réglée.

   Toutes les prestations proposées par le Centre Équestre peuvent être payées \textit{via} les moyens suivants :
   \begin{itemize}
      \item Espèces de toute valeur (la direction se réserve le droit de refuser les coupures de 100, 200 et 500 euros) ;
      \item chèque en euros (à l'ordre de \texttt{SARL Centre Équestre de Sauveterre}) ;
      \item virement bancaire ;
      \item carte bancaire ;
      \item chèques vacances ;
      \item coupons sport ;
      \item carte cadeau émise par le Centre Équestre, en cours de validité.
   \end{itemize}

   La direction se réserve le droit de refuser tout moyen de paiement ne figurant pas dans cette liste.

   Les \og cartes cadeau \fg{} sont des coupons prépayés émis par le Centre Équestre pouvant être déclinés pour tout type de prestation, et édités après paiement de la prestation voulue ou à titre promotionnel.
   Elles ont vocation à être offertes, et sont d'une validité variable en fonction de la prestation.
   Une carte cadeau ne peut être ni échangée, ni remboursée, ni voir sa validité prolongée (sauf cas de force majeure, sur accord de la direction).
   Elle devra être présentée le jour de l'activité et constitue à ce titre un moyen de paiement.

   Pour les cartes de cours et forfaits, un paiement en plusieurs fois est possible en accord avec le gérant.

   \paragraph[Modalités de remboursement :  ]{Modalités de remboursement :\\}
   Concernant les modalités de remboursement des balades, se référer aux \textit{Conditions Générales applicables aux Balades}.

   Pour toutes les autres prestations, le remboursement est décidé au cas par cas entre la direction et les intéressés.

   \subsection{Propriétaires de chevaux}\label{subsec:proprietaires-de-chevaux}
   Les propriétaires auront pris connaissance des \textit{Conditions Générales de prise en pension}, présentées en annexe.

   \subsection{Armoires de rangement}\label{subsec:armoires-de-rangement}
   Afin de ranger et protéger les affaires, des armoires de rangement fermées à clé ou cadenas peuvent être louées, deux tailles étant proposées.
   Ces casiers peuvent être loués par tout cavalier du club, propriétaire ou non, seul ou à plusieurs, dans la limite des disponibilités.
   De plus, tout locataire de casier devra indiquer son nom sur son casier, en utilisant l'étiquette prévue à cet effet.

   Le paiement est mensuel, et \emph{tout casier non réglé pendant une période de 6 mois sera libéré et les affaires rangées à l'accueil jusqu'à régularisation du paiement}.
   Évidemment, tout casier libéré peut être attribué à un autre cavalier en faisant la demande.

   Sont mis à disposition des cavaliers, dans chaque casier, \textit{a minima} un porte-selle et un porte-filet pour les petits casiers, et deux de chaque dans les grands.
   Les cavaliers sont libres de ranger ce qu'ils veulent dans leurs casiers, cependant \emph{merci de ne pas y ranger de nourriture, ou alors la garder dans des boîtes hermétiques} pour des raisons évidentes d'hygiène.

   \subsection{Autorisation de reproduction d'images et cession de droits}\label{subsec:autorisation-de-reproduction-d'images-et-cession-de-droits}
   Conformément aux dispositions relatives au droit à l'image et au droit d'auteur, les adhérents autorisent le Centre Équestre ainsi que son Association à capter, utiliser ou diffuser des images photo ou vidéo prises dans le cadre des activités du Club, sur les supports suivants :
   \begin{itemize}
      \item
      site internet du Club ;
      \item
      journal d'information du Club ;
      \item
      documents de présentation et promotion des activités de la SARL Centre Équestre de Sauveterre ;
      \item
      journaux de la presse quotidienne locale ;
      \item
      journaux spécialisés des sports équestres.
   \end{itemize}
   Dans les conditions d'utilisation définies ci-dessus, les adhérents renoncent à tout recours ou toute indemnité contre le Centre Équestre ou son Association.

   Tout refus devra être notifié par écrit à la direction.


   \vspace{4cm}

   Fait aux Sables d'Olonne, le \today

   \vspace{2cm}

   \begin{flushright}
      \rule{6cm}{0.5pt}

      {Stéphane MARAIS, Directrice \kern 9pt}

   \end{flushright}

\end{document}
